\section*{Chapter \ref{ch:b3:panel-data} --- Panel Data}

\begin{solution}{exo:b3:panel-data:1}
\begin{enumerate}
\item Between: 2.7 tonnes per log point, comparing countries' averages with each
other. Within: 1.0, comparing each country's years with its own average.
% figdata: co2_stats
\item An endowment of coal, oil or gas, which makes energy cheap and exports
valuable; a cold climate, which raises both heating and, through the
industries it favours, income. A large car-based territory also fits.
\item Most of the variation in income is between countries: incomes differ by a
factor of a hundred across countries, but one country's income moves by a few
tenths of a log point around its own average over three decades.
\end{enumerate}
\end{solution}

\begin{solution}{exo:b3:panel-data:2}
\begin{enumerate}
\item $\bar x = 3.5$ and $\bar y = 7$. The cross-products sum to $12.5 + 4.5 + 1
+ 1 + 4.5 + 12.5 = 36$ and the squares of $x - \bar x$ to $17.5$: the pooled
slope is $36/17.5 = 2.057$.
\item Unit A: $\bar x_A = 2$, $\bar y_A = 11/3$; deviations $\ddot x = (-1, 0,
1)$, $\ddot y = (-5/3, 1/3, 4/3)$. Unit B: $\bar x_B = 5$, $\bar y_B = 31/3$;
$\ddot x = (-1, 0, 1)$, $\ddot y = (-4/3, -1/3, 5/3)$. Then $\sum \ddot x \ddot
y = 5/3 + 4/3 + 4/3 + 5/3 = 6$ and $\sum \ddot x^2 = 4$: the within slope is
$1.5$.
\item The two clouds each rise by about 1.5 per unit of $x$, but B's lies
higher and further right, so a single line through all six points (2.057) is
steeper than either.
\end{enumerate}
\end{solution}

\begin{solution}{exo:b3:panel-data:3}
Unit A: $\Delta x = (1, 1)$, $\Delta y = (2, 1)$; unit B: $\Delta x = (1, 1)$,
$\Delta y = (1, 2)$. The slope without a constant is $\sum \Delta x \Delta y /
\sum \Delta x^2 = (2 + 1 + 1 + 2)/4 = 1.5$. Every $\Delta x$ equals 1, so
$\Delta x$ is collinear with a constant: a regression with a constant cannot
separate the slope from the intercept.
\end{solution}

\begin{solution}{exo:b3:panel-data:4}
First, split the panel into 1990--2005 and 2006--2022 and estimate the within
slope in each; with constant unit effects and a constant slope, the two agree up
to sampling error. Second, compare the within and first-difference estimates: if
a slowly changing unit effect (a drift in each country's energy mix) were left
in the error, the two would weight it differently and disagree. A third check:
the residuals of the within regression should show no trend inside countries.
\end{solution}

\begin{solution}{exo:b3:panel-data:5}
The tested solution is \code{ex05_year_effects.py}, in the chapter's
\code{exercises/} folder.
\begin{enumerate}
\item Country effects: slope 0.999, conventional s.e.\ 0.122, clustered 0.265,
ratio 2.17. Country and year effects: 2.068, 0.164 and 0.450, ratio 2.74.
% output: ex05_year_effects
\item The cross-country mean of log income rose by 0.576. A common rise in
income, with emissions per person in several large rich countries falling over
the same years, makes the year-to-year co-movement shared by all countries
negative or weak; the year effects remove it, and what remains, each country's
movement relative to the world's that year, has a steeper slope.
\end{enumerate}
\end{solution}

\begin{solution}{exo:b3:panel-data:6}
\begin{enumerate}
\item Downward: the omitted unit effect is negatively correlated with income, so
the between slope is below the true slope.
\item No. On this story alone, the between slope would be smaller than an
unbiased within slope, yet the chart shows it larger. Some fixed
characteristics that raise income must also raise emissions (fossil
endowments, climate) and outweigh the clean-technology effect; or the within
slope is itself small for other reasons, such as measurement error or the
short-run nature of within variation.
\end{enumerate}
\end{solution}

\begin{solution}{exo:b3:panel-data:7}
\begin{enumerate}
\item $(0.265/0.122)^2 = 2.172^2 = 4.72$.
\item $6{,}419/4.72 \approx 1{,}360$ effective observations.
\item Conventional: $0.999 \pm 1.96 \times 0.122 = 0.999 \pm 0.239$, from
0.760 to 1.238. Clustered: $0.999 \pm 0.519$, from 0.480 to 1.518.
\end{enumerate}
\end{solution}

\begin{solution}{exo:b3:panel-data:8}
\begin{enumerate}
\item Random effects moves; fixed effects does not. The within transform
removes $a_i$ entirely, so its correlation with $x$ cannot reach the estimate.
Random effects removes only the share $\theta$ of each unit's mean and leaves
$(1 - \theta)a_i$ in the error (\cref{prop:b3:panel-data:re}).
\item $(1 - 0.610)^2 = 0.152$. The true value is $1/(1 + 6) = 0.143$, so the
true $\theta$ is $1 - \sqrt{0.143} = 0.622$; the estimate, 0.610, is close.
\end{enumerate}
\end{solution}

\begin{solution}{exo:b3:panel-data:9}
The tested solution is \code{ex09_fertility.py}, in the chapter's
\code{exercises/} folder.
\begin{enumerate}
\item Over 213 countries: between $-0.778$, within $-0.837$ births per woman per
log point of income. % output: ex09_fertility
\item The fixed differences between countries are not strongly correlated with
income in a way that also moves fertility; whatever the fixed effects absorb
does not change the slope much. The two slopes can still both be far from a
causal effect: income and fertility move together inside countries for reasons,
such as women's schooling, that change over time.
\end{enumerate}
\end{solution}

\begin{solution}{exo:b3:panel-data:10}
With $T = 2$, $\bar x_i = (x_{i1} + x_{i2})/2$, so $\ddot x_{i1} = -\Delta x_i/2$
and $\ddot x_{i2} = \Delta x_i/2$, with $\Delta x_i = x_{i2} - x_{i1}$; the same
holds for $y$. Then
\[
  \sum_i\sum_t \ddot x_{it}\ddot y_{it} = \sum_i 2 \cdot \frac{\Delta x_i \Delta y_i}{4} = \frac12 \sum_i \Delta x_i\Delta y_i,
  \qquad
  \sum_i\sum_t \ddot x_{it}^2 = \frac12 \sum_i \Delta x_i^2,
\]
and the ratio is $\sum_i \Delta x_i \Delta y_i / \sum_i \Delta x_i^2$, the
\omterm{def:b3:panel-data:first-difference}{first-difference estimator} without a constant.
\end{solution}

\begin{solution}{exo:b3:panel-data:11}
The tested solution is \code{ex11_nickell.py}, in the chapter's
\code{exercises/} folder.
\begin{enumerate}
\item Mean estimates 0.168, 0.337 and 0.421 for $T = 5$, 10 and 20; biases
$-0.332$, $-0.163$ and $-0.079$. % output: ex11_nickell
\item The approximation gives $-0.375$, $-0.167$ and $-0.079$. It is good from
$T = 10$ on and overstates the bias at $T = 5$, as an approximation for large
$T$ may.
\end{enumerate}
\end{solution}

\begin{solution}{exo:b3:panel-data:12}
\begin{enumerate}
\item $\ddot u_{it} = u_{it} - \bar u_i$ contains $u_{is}$ for every $s$,
including $s < t$. Under feedback, $x_{it}$ depends on $u_{i,t-1}$, so
$\Cov(x_{it}, \bar u_i) \neq 0$ and hence $\Cov(\ddot x_{it}, \ddot u_{it}) \neq
0$, even if $\Cov(x_{it}, u_{it}) = 0$.
\item Under strict exogeneity, $x_{i,t+1}$ is unrelated to $u_{it}$ and its
coefficient is zero. If a positive shock this year raises $x$ next year, the
lead is positively correlated with $u_{it}$: its coefficient is positive.
\end{enumerate}
\end{solution}

\begin{solution}{pb:b3:panel-data:firms}
\textbf{1.} Across firm-years, 1~\% more capital goes with 0.87~\% more output.
% figdata: firms_results

\textbf{2.} $0.30 + 0.5 \times 0.9512/0.8675 = 0.30 + 0.5 \times 1.0965 = 0.30
+ 0.548 = 0.848$.

\textbf{3.} The pooled estimate is 0.868, close to the prediction; the rest
comes from the subsidy and the year shocks, also left in the pooled error.

\textbf{4.} It is not recorded. Any measure of it --- a survey score, a
manager's degree --- would be partial and noisy, and the omitted part would
remain.

\textbf{5.} Downward: the covariance of capital with management would be
negative.

\textbf{6.} $0.0094/0.0045 = 2.09$. No: the standard error measures sampling
noise around the wrong number. A biased estimate can be very precise.

\textbf{7.} Within a firm, a year in which capital is 1~\% above the firm's own
average has output 0.35~\% above its average. The comparisons are between
years of the same firm.

\textbf{8.} The errors are serially correlated, and the two estimators weight
the years differently; and the year shocks common to all firms enter both
differently, because capital trends upward and the within transform compares
each year with the firm's whole-period average, while differences compare
adjacent years. Measurement error in capital would also hit them differently.

\textbf{9.} Shocks common to all firms in a year (demand, prices). Capital
rises on average over time, so any common year shock that happens to line up
with that trend is attributed to capital by the one-way within estimate.

\textbf{10.} $0.3182 \pm 1.96 \times 0.0147 = 0.3182 \pm 0.0288$, from 0.289 to
0.347.

\textbf{11.} The movements of each firm's capital around its own average, net of
the movement common to all firms in that year.

\textbf{12.} No. The advice changes over time (it starts in 2015) and moves both
capital and output; it is a time-varying confounder, which neither firm nor year
effects remove.

\textbf{13.} The panel picked better-run firms, and management raises output.
Capital in the regression captures only part of management, so the subsidy
dummy picks up the rest.

\textbf{14.} Pooled 0.122, two-way 0.110. The pooled estimate is larger, as
Question~13 predicts, but here by little, because capital already absorbs much
of the management difference.

\textbf{15.} $\eu^{0.1101} - 1 = 0.116$: output 11.6~\% higher.

\textbf{16.} The change in output of subsidised firms from before 2018 to after,
compared with the change of unsubsidised firms over the same years.

\textbf{17.} Without the subsidy, the output of treated and untreated firms would
have moved in parallel after 2018: their gap could be anything, but not its
change.

\textbf{18.} Estimate the gap between treated and untreated firms year by year
for 2011--2017; it should be flat. A trend in it before 2018 would undermine the
assumption.

\textbf{19.} Upward: output that has just fallen tends to recover by itself, and
the recovery would be counted as the subsidy's effect.

\textbf{20.} $0.0181/0.0110 = 1.645$; its square is $2.71$.

\textbf{21.} iid: $0.1101 \pm 1.96 \times 0.0110 = 0.1101 \pm 0.0216$, from
0.0885 to 0.1317. Clustered: $0.1101 \pm 0.0355$, from 0.0746 to 0.1456.

\textbf{22.} The errors persist within a firm from year to year, and the
regressor persists too: the subsidy is on for eight years in a row for the same
firms. Together they make the firm-years of one firm far from independent.

\textbf{23.} $2.71 = 1 + 14\rho$, so $\rho = 1.71/14 = 0.122$.

\textbf{24.} $4{,}500/2.71 \approx 1{,}660$ firm-years.

\textbf{25.} Inside firms, net of year shocks, the elasticity of output with
respect to capital is 0.318, with a clustered 95~\% interval of 0.289 to 0.347.
The training subsidy raised output by 0.110 log points (11.6~\%), with a
clustered 95~\% interval of 0.075 to 0.146, if treated and untreated firms would
otherwise have moved in parallel.
\end{solution}
